\textbf{Phần 1. Trắc nghiệm lựa chọn}
\setcounterpageref{ex}{0}

\begin{ex} %[3H3N22]
Một đoạn thẳng $AB$ dài $12\text{ cm}$. Điểm $M$ là trung điểm của $AB$ thì $AM$ bằng
\choice
{$4\text{ cm}$}
{$8\text{ cm}$}
{\True $6\text{ cm}$}
{$12\text{ cm}$}
\loigiai{
$AM = 12 : 2 = 6\text{ cm}$.\\
\textbf{Đáp án của câu hỏi là C}
}
\end{ex}

\begin{ex} %[3H3N22]
Hình tròn tâm $O$ có bán kính $r = 5\text{ cm}$. Đường kính của hình tròn là
\choice
{$5\text{ cm}$}
{$15\text{ cm}$}
{$20\text{ cm}$}
{\True $10\text{ cm}$}
\loigiai{
$5 \times 2 = 10\text{ cm}$.\\
\textbf{Đáp án của câu hỏi là D}
}
\end{ex}

\begin{ex} %[3H3H22]
Hình vuông có cạnh dài $6\text{ cm}$. Chu vi hình vuông đó là
\choice
{$18\text{ cm}$}
{\True $24\text{ cm}$}
{$30\text{ cm}$}
{$36\text{ cm}$}
\loigiai{
$6 \times 4 = 24\text{ cm}$.\\
\textbf{Đáp án của câu hỏi là B}
}
\end{ex}

\begin{ex} %[3H3H22]
Hình chữ nhật có chiều dài $10\text{ cm}$, chiều rộng $6\text{ cm}$. Chu vi hình chữ nhật đó là
\choice
{$16\text{ cm}$}
{$24\text{ cm}$}
{$60\text{ cm}$}
{\True $32\text{ cm}$}
\loigiai{
$(10 + 6) \times 2 = 32\text{ cm}$.\\
\textbf{Đáp án của câu hỏi là D}
}
\end{ex}

\begin{ex} %[3H3H22]
Khối lập phương có bao nhiêu đỉnh và bao nhiêu cạnh?
\choice
{$6$ đỉnh, $8$ cạnh}
{$12$ đỉnh, $8$ cạnh}
{\True $8$ đỉnh, $12$ cạnh}
{$6$ đỉnh, $12$ cạnh}
\loigiai{
Khối lập phương có 8 đỉnh và 12 cạnh.\\
\textbf{Đáp án của câu hỏi là C}
}
\end{ex}

\begin{ex} %[3H3H22]
Khối hộp chữ nhật có bao nhiêu mặt?
\choice
{\True $6$ mặt}
{$4$ mặt}
{$8$ mặt}
{$12$ mặt}
\loigiai{
Khối hộp chữ nhật có 6 mặt.\\
\textbf{Đáp án của câu hỏi là A}
}
\end{ex}

\begin{ex} %[3H3V22]
Một tam giác có độ dài các cạnh là $6\text{ cm}, 8\text{ cm}, 10\text{ cm}$. Chu vi hình tam giác là
\choice
{$22\text{ cm}$}
{$26\text{ cm}$}
{\True $24\text{ cm}$}
{$28\text{ cm}$}
\loigiai{
$6 + 8 + 10 = 24\text{ cm}$.\\
\textbf{Đáp án của câu hỏi là C}
}
\end{ex}

\begin{ex} %[3H3V22]
Một hình tròn có đường kính $14\text{ cm}$. Bán kính của nó là
\choice
{\True $7\text{ cm}$}
{$5\text{ cm}$}
{$8\text{ cm}$}
{$28\text{ cm}$}
\loigiai{
$14 : 2 = 7\text{ cm}$.\\
\textbf{Đáp án của câu hỏi là A}
}
\end{ex}

\begin{ex} %[3H3V22]
Một hình vuông có chu vi $36\text{ cm}$. Cạnh của hình vuông đó là
\choice
{$8\text{ cm}$}
{$10\text{ cm}$}
{$12\text{ cm}$}
{\True $9\text{ cm}$}
\loigiai{
$36 : 4 = 9\text{ cm}$.\\
\textbf{Đáp án của câu hỏi là D}
}
\end{ex}

\begin{ex} %[3H3V22]
Đoạn thẳng $AB$ có trung điểm $M$. Biết $AM = 7\text{ cm}$. Độ dài đoạn thẳng $AB$ là
\choice
{$12\text{ cm}$}
{\True $14\text{ cm}$}
{$16\text{ cm}$}
{$21\text{ cm}$}
\loigiai{
$7 \times 2 = 14\text{ cm}$.\\
\textbf{Đáp án của câu hỏi là B}
}
\end{ex}

\begin{ex} %[3H3V22]
Hình chữ nhật có nửa chu vi là $20\text{ cm}$, chiều rộng là $7\text{ cm}$. Chiều dài hình chữ nhật là
\choice
{\True $13\text{ cm}$}
{$12\text{ cm}$}
{$14\text{ cm}$}
{$15\text{ cm}$}
\loigiai{
$20 - 7 = 13\text{ cm}$.\\
\textbf{Đáp án của câu hỏi là A}
}
\end{ex}

\begin{ex} %[3H3C22]
Một khối lập phương bằng gỗ có cạnh $4\text{ cm}$. Người ta sơn đỏ tất cả 6 mặt của khối gỗ. Tổng diện tích các mặt được sơn là (biết diện tích mỗi mặt là $4 \times 4 = 16\text{ cm}^2$):
\choice
{$64\text{ cm}^2$}
{\True $96\text{ cm}^2$}
{$80\text{ cm}^2$}
{$100\text{ cm}^2$}
\loigiai{
Diện tích 6 mặt: $16 \times 6 = 96\text{ cm}^2$.\\
\textbf{Đáp án của câu hỏi là B}
}
\end{ex}

\textbf{Phần 2. Trắc nghiệm đúng sai}
\setcounterpageref{ex}{0}

\begin{ex} %[3H3H22]
Xét tính đúng sai về hình học:
\choiceTF
{\True Hình tròn có vô số bán kính và tất cả đều bằng nhau.}
{\True Hình chữ nhật có 4 góc vuông.}
{Khối lập phương có 12 đỉnh và 8 cạnh.}
{\True Khối lập phương có 8 đỉnh và 12 cạnh.}
\loigiai{
\begin{itemize}
    \item Hình tròn có vô số bán kính bằng nhau: Đúng.
    \item Hình chữ nhật có 4 góc vuông: Đúng.
    \item Có 8 đỉnh và 12 cạnh (không phải 12 đỉnh 8 cạnh). Ý 3 Sai, Ý 4 Đúng.
\end{itemize}
\begin{itemize}
    \item a) \textbf{Đúng}.
    \item b) \textbf{Đúng}.
    \item c) \textbf{Sai}.
    \item d) \textbf{Đúng}.
\end{itemize}
}
\end{ex}

\begin{ex} %[3H3V22]
Cho đoạn thẳng $CD = 24\text{ cm}$ có $I$ là trung điểm. Xét tính đúng sai:
\choiceTF
{\True Đoạn thẳng $CI$ dài $12\text{ cm}$.}
{\True Điểm $I$ nằm chính giữa đoạn thẳng $CD$.}
{Nếu $K$ là trung điểm của $CI$ thì $CK = 8\text{ cm}$.}
{\True Nếu $K$ là trung điểm của $CI$ thì $CK = 6\text{ cm}$.}
\loigiai{
\begin{itemize}
    \item $CI = 24 : 2 = 12\text{ cm}$. Đúng.
    \item $I$ là trung điểm nên nằm chính giữa: Đúng.
    \item $CK = 12 : 2 = 6\text{ cm} \neq 8\text{ cm}$. Ý 3 Sai, Ý 4 Đúng.
\end{itemize}
\begin{itemize}
    \item a) \textbf{Đúng}.
    \item b) \textbf{Đúng}.
    \item c) \textbf{Sai}.
    \item d) \textbf{Đúng}.
\end{itemize}
}
\end{ex}

\textbf{Phần 3. Trả lời ngắn}
\setcounterpageref{ex}{0}

\begin{ex} %[3H3H22]
Một hình tròn có bán kính $8\text{ cm}$. Tính đường kính hình tròn theo xăng-ti-mét.
\shortans{16}
\loigiai{
\textbf{Đáp số:} $16$\\
$8 \times 2 = 16\text{ cm}$.
}
\end{ex}

\begin{ex} %[3H3H22]
Tính chu vi hình vuông có cạnh dài $6\text{ cm}$.
\shortans{24}
\loigiai{
\textbf{Đáp số:} $24$\\
$6 \times 4 = 24\text{ cm}$.
}
\end{ex}

\begin{ex} %[3H3V22]
Tính chu vi hình chữ nhật có chiều dài $14\text{ cm}$, chiều rộng $6\text{ cm}$.
\shortans{40}
\loigiai{
\textbf{Đáp số:} $40$\\
$(14 + 6) \times 2 = 40\text{ cm}$.
}
\end{ex}

\begin{ex} %[3H3V22]
Đoạn thẳng $AB$ có trung điểm $M$ và $AM = 9\text{ cm}$. Tính độ dài đoạn thẳng $AB$ theo xăng-ti-mét.
\shortans{18}
\loigiai{
\textbf{Đáp số:} $18$\\
$9 \times 2 = 18\text{ cm}$.
}
\end{ex}

\begin{ex} %[3H3V22]
Một hình vuông có chu vi bằng chu vi hình chữ nhật có chiều dài $10\text{ cm}$, chiều rộng $6\text{ cm}$. Tính cạnh của hình vuông đó theo xăng-ti-mét.
\shortans{8}
\loigiai{
\textbf{Đáp số:} $8$\\
Chu vi: $(10 + 6) \times 2 = 32\text{ cm}$. Cạnh hình vuông: $32 : 4 = 8\text{ cm}$.
}
\end{ex}

\begin{ex} %[3H3C22]
Một khung dây thép hình lập phương có cạnh dài $8\text{ cm}$. Tính tổng chiều dài dây thép để làm khung (12 cạnh) theo xăng-ti-mét.
\shortans{96}
\loigiai{
\textbf{Đáp số:} $96$\\
$8 \times 12 = 96\text{ cm}$.
}
\end{ex}
